\documentclass[11pt,a4paper]{article}
\usepackage[T2A]{fontenc}
\usepackage[utf8]{inputenc}
\usepackage[russian]{babel}
\usepackage{geometry}
\usepackage{amsmath,amssymb,mathtools}
\usepackage{array,booktabs}
\usepackage{xcolor}
\usepackage{tikz}
\usepackage[unicode,colorlinks=true,linkcolor=blue!55!black]{hyperref}
\geometry{left=19mm,right=19mm,top=20mm,bottom=21mm}
\setlength{\parindent}{0pt}
\setlength{\parskip}{5pt}
\setlength{\emergencystretch}{2em}
\renewcommand{\arraystretch}{1.25}
\hypersetup{pdftitle={Ревью домашней работы — 14 октября 2026 года},pdfauthor={Лёвин Артём Александрович}}
\begin{document}

% Титульный лист.
\thispagestyle{empty}
\begin{center}
\vspace*{0.31\textheight}
{\Huge\bfseries Ревью домашней работы\par}
\vspace{20mm}
{\large Автор: Лёвин Артём Александрович\par}
\vspace{8mm}
{\large 14 октября 2026 года\par}
\end{center}
\vfill
\begin{center}
\begin{tikzpicture}[x=1cm,y=1cm]
\draw[blue!35!black,line width=0.75pt]
  (0,0) .. controls (2.4,0.50) and (4.5,-0.48) .. (7.0,0);
\end{tikzpicture}
\end{center}
\newpage

\section*{Сводная таблица ошибок}
В таблицу внесены только задания, при выполнении которых обнаружены ошибки в ходе решения.

{\small
\begin{center}
\begin{tabular}{@{}>{\centering\arraybackslash}p{14mm}
                >{\raggedright\arraybackslash}p{29mm}
                >{\raggedright\arraybackslash}p{64mm}
                >{\centering\arraybackslash}p{24mm}
                >{\centering\arraybackslash}p{24mm}@{}}
\toprule
\textbf{№} & \textbf{Тема} & \textbf{Суть ошибки} & \textbf{Ваш ответ} & \textbf{Правильный ответ}\\
\midrule
\hyperref[task:10]{10} & Вычитание смешанных чисел со скобками &
При заимствовании одной целой части получено $111/72$ вместо $114/72$; следующий результат вычитания не следует из записанного выражения. &
$3\dfrac{49}{72}$ & $3\dfrac{49}{72}$\\
\bottomrule
\end{tabular}
\end{center}}

\newpage
\thispagestyle{empty}
\vspace*{\fill}
\begin{center}
{\LARGE\bfseries Разбор ошибок}
\end{center}
\vspace*{\fill}
\newpage

\phantomsection\label{task:10}
\section*{Задание № 10}

\textbf{Текст задания.} Выполните вычитание:
\[
9\frac{7}{12}-\left(4\frac{5}{18}+1\frac{5}{8}\right).
\]

\textbf{Ваше решение.} Сначала Вы верно привели дробные части слагаемых в скобках к знаменателю $72$ и сложили их:
\[
4\frac{20}{72}+1\frac{45}{72}=5\frac{65}{72}.
\]
Затем записали следующие преобразования:
\[
\begin{aligned}
9\frac{7}{12}-5\frac{65}{72}
 &=9\frac{42}{72}-5\frac{65}{72}\\
 &=8\frac{111}{72}-5\frac{65}{72}\\
 &=3\frac{49}{72}.
\end{aligned}
\]

\textbf{Ваш ответ:} $3\dfrac{49}{72}$.

\textbf{\textcolor{red}{Ошибка:}} неверное преобразование смешанного числа при заимствовании одной целой части; последующее вычитание не согласуется с записанным числителем.

\textbf{Где именно возникает ошибка.}

Последний верный шаг:
\[
9\frac{42}{72}-5\frac{65}{72}.
\]
Первый неверный шаг --- замена
\[
9\frac{42}{72}\quad\text{на}\quad
\textcolor{red}{8\frac{111}{72}}.
\]
При заимствовании одной целой части правильно получить $8\dfrac{114}{72}$, а не $8\dfrac{111}{72}$.

\textbf{Почему этот шаг неверен.}

Смешанное число $9\dfrac{42}{72}$ означает сумму девяти целых и сорока двух семьдесят вторых. Для вычитания $65/72$ дробной части $42/72$ недостаточно, поскольку $42<65$. Поэтому одну из девяти целых частей нужно представить в виде дроби со знаменателем $72$:
\[
1=\frac{72}{72}.
\]
После этого останется восемь целых, а к дробной части прибавятся именно семьдесят две семьдесят вторых:
\[
9\frac{42}{72}
=8+\frac{72}{72}+\frac{42}{72}
=8+\frac{114}{72}
=8\frac{114}{72}.
\]
Значит, числитель должен равняться $72+42=114$. Число $111$ из такого преобразования не получается. Следующий переход также некорректен: если формально вычесть $65$ из ошибочно записанного $111$, получится $46$, а не $49$. Поэтому совпадение Вашего окончательного ответа с верным не делает промежуточную цепочку равенств правильной.

\textbf{\textcolor{blue!60!black}{Ключевая идея:}} если дробная часть уменьшаемого меньше дробной части вычитаемого, нужно занять одну целую часть. При знаменателе $72$ одна целая равна $72/72$, поэтому к числителю дробной части прибавляется ровно $72$.

\textbf{Исправление от последнего правильного шага.}
\[
\begin{aligned}
9\frac{42}{72}-5\frac{65}{72}
 &=8\frac{114}{72}-5\frac{65}{72}\\
 &=(8-5)+\frac{114-65}{72}\\
 &=3\frac{49}{72}.
\end{aligned}
\]

\textbf{Полное правильное решение.}

\emph{Шаг 1. Выполним сложение в скобках.} Наименьший общий знаменатель для дробей $5/18$ и $5/8$ равен $72$, так как $72:18=4$ и $72:8=9$. Получаем:
\[
\frac{5}{18}=\frac{20}{72},\qquad
\frac{5}{8}=\frac{45}{72}.
\]
Складываем целые и дробные части отдельно:
\[
4\frac{5}{18}+1\frac{5}{8}
=4\frac{20}{72}+1\frac{45}{72}
=5\frac{65}{72}.
\]

\emph{Шаг 2. Приведём дробную часть уменьшаемого к знаменателю $72$.} Поскольку $72:12=6$, имеем:
\[
\frac{7}{12}=\frac{7\cdot6}{12\cdot6}=\frac{42}{72}.
\]
Поэтому исходный пример превращается в
\[
9\frac{42}{72}-5\frac{65}{72}.
\]

\emph{Шаг 3. Займём одну целую часть.} Так как $42/72<65/72$, представим девять целых как восемь целых и ещё одну целую, то есть $72/72$:
\[
9\frac{42}{72}=8\frac{72+42}{72}=8\frac{114}{72}.
\]

\emph{Шаг 4. Вычтем целые и дробные части.}
\[
8\frac{114}{72}-5\frac{65}{72}
=(8-5)+\frac{114-65}{72}
=3\frac{49}{72}.
\]
Дробь $49/72$ сокращать нельзя, поскольку числитель и знаменатель не имеют общего делителя больше единицы.

\textbf{Проверка.} Прибавим к полученной разности вычитаемое:
\[
\begin{aligned}
3\frac{49}{72}+5\frac{65}{72}
 &=8\frac{114}{72}
 =9\frac{42}{72}
 =9\frac{7}{12}.
\end{aligned}
\]
Получено исходное уменьшаемое, значит, значение разности найдено верно.

\textbf{\textcolor{teal!60!black}{Итог:}} правильный ответ --- $3\dfrac{49}{72}$. Он совпал с Вашим, но при заимствовании единицы необходимо получить $114/72$, после чего вычитание $114-65=49$ становится обоснованным.

\medskip
\textbf{Задача на закрепление.}
Вычислите:
\[
8\frac{5}{12}-\left(3\frac{5}{9}+1\frac{1}{8}\right).
\]

\vfill
\begin{center}
{\footnotesize Лёвин Артём Александрович эксклюзивно для Ксении Васильченко}
\end{center}

\end{document}
